\chapter{Geometria di contatto}

In questo capitolo tratteremo la \emph{geometria di contatto}: essa rappresenta il cardine formale di tutta la trattazione dell'elaborato \cite{jiang_thermodynamics}.
Daremo delle definizioni per inquadrarne la struttura rispetto alle nozioni descritte nei capitoli precedenti; si mostrerà la relazione con la sua «sorella», la \emph{geometria simplettica}, definita per varietà a dimensioni pari; si cercherà di completare la descrizione di esempi che rendano più intuitivi i concetti: si vedrà come la struttura di geometria di contatto ha applicazioni in analisi matematica e meccanica hamiltoniana.

\section{Forme di contatto}

\begin{definition}
    Una \textbf{distribuzione D di Frobenius} \cite{pirola_geometria} di rango k è un \emph{sottofibrato} di rango k di $T\mathcal{M}$.
    Un campo vettoriale $X \in \mathcal{X}(\mathcal{M})$ appartiene alla distribuzione $D$ se per ogni $p \in \mathcal{M}$, $X(p) \in D_p$. In altre parole $X$ è una sezione di D.
\end{definition}
\[
    X \in D \quad \mathrm{se} \quad \forall p \in \mathcal{M}, \; \;  X(p) \in D_p
\]

\begin{definition}
    Sia $\mathcal{M}$ una varietà liscia $(2n + 1)$-dimensionale, una \textbf{struttura di contatto} è una distribuzione liscia $\xi$ di iperpiani, $\xi \subset T \mathcal{M}$ tale che qualsiasi 1-forma differenziale $\alpha$ che definisce $\xi$ soddisfa la condizione di non-integrabilità massimale.

    \begin{equation}
        \alpha \wedge (\odif{\alpha})^n \neq 0 \tag*{\textbf{Condizione di non-integrabilità massimale}}
    \end{equation}

    dove $ (\odif{\alpha})^n = \bigwedge^n \odif{\alpha}$.
    Una 1-forma $\alpha$ \emph{globale} è detta \textbf{forma di contatto} e $(\mathcal{M}, \xi)$ è una \textbf{varietà di contatto}.
\end{definition}

Prestiamo attenzione a due aspetti. Con $\alpha$ che definisce la distribuzione di iperpiani $\xi$ si intende che per ogni $p \in \mathcal{M}$, l'iperpiano $\xi$ è composto esclusivamente da tutti i vettori tangenti tali che $\alpha (X) = 0$, cioè $\xi = \ker \alpha$.
Da questo si vede che se $\mathcal{N}$ è una sottovarietà tangente a $\xi$, allora \( \alpha^* : T^*(\mathcal{M}) \rightarrow T^*(\mathcal{M})\)

Inoltre, una forma di contatto $\alpha$ non è univocamente determinata: possiamo moltiplicare una qualsiasi $\alpha$ per una funzione liscia non-nulla e si ottiene una forma di contatto.
Ne diamo una piccola ``dimostrazione'':
\begin{enumerate}
    \item $\ker \alpha = \ker (f\alpha)$
    \item Calcoliamo
          \[
              (f\alpha) \wedge \bigl(\odif{(f\alpha)}\bigr)^n
              = f\alpha \wedge \bigl(f\,\odif{\alpha} + \odif{f}\wedge\alpha\bigr)^n .
          \]
          Usiamo il teorema binomiale, $(a+b)^n = \sum_{k=0}^n \binom{n}{k} a^{n-k}b^k$
          (il che è lecito perché $f\,\odif{\alpha}$ e $\odif{f}\wedge\alpha$ sono $2$-forme, quindi commutano):
          \begin{align*}
              f\alpha \wedge \bigl(f\,\odif{\alpha} + \odif{f}\wedge\alpha\bigr)^n
               & = f\alpha \wedge \Bigl[
                  f^n \odif{\alpha}^n
                  + \binom{n}{1} f^{n-1} (\odif{\alpha})^{n-1} \wedge \odif{f}\wedge\alpha
              + \dots                                                                                                                  \\
               & \hspace{4em} \dots
                  + \binom{n}{n-1} f\,\odif{\alpha} \wedge (\odif{f}\wedge\alpha)^{n-1}
                  + (\odif{f}\wedge\alpha)^n
              \Bigr] =                                                                                                                 \\
               & = f \alpha \wedge (f \odif{\alpha})^n + f \alpha \wedge \left[
                  n (f \odif{\alpha})^{n-1} \wedge \odif{f} \wedge \alpha
              \right] =                                                                                                                \\
               & = f^{n+1} \alpha \wedge (\odif{\alpha})^n - f n (f \odif{\alpha})^{n-1} \wedge \odif{f} \wedge \alpha \wedge \alpha = \\
               & = f^{n+1}\,  \alpha \wedge \odif{\alpha}^n
          \end{align*}
          dove è stato usato \(\alpha \wedge \alpha = 0\) per la sua antisimmetria.
          Allora otteniamo esattamente
          \begin{equation}
              f \alpha \wedge (\odif{f \alpha})^n = f^{n+1} \alpha \wedge (\odif{\alpha})^n
          \end{equation}
\end{enumerate}

\begin{lemma}
    Sia $(\mathcal{M}, \xi)$ una varietà di contatto, supponiamo \(\alpha_1, \alpha_2\) forme di contatto della distribuzione di iperpiani $\xi$, allora
    \[
        \alpha_1 = f \alpha_2
    \]
    per qualche \(f : \mathcal{M} \longrightarrow \mathbb{R}\) liscia non-nulla.
\end{lemma}

\begin{proof}
    Poiché \(\alpha_1, \alpha_2\) sono forme di contatto per $\xi$, vale che $\ker(\alpha_1) = \xi = \ker(\alpha_2)$.
    $\alpha_1, \alpha_2$ possono essere viste come covettori nello spazio cotangente che sono uno multiplo dell'altro tramite prodotto per scalare.
    Costruiamo $f$: sia $p \in \mathcal{M}$, scegliamo un generico vettore non-nullo $v_p \in T_p \mathcal{M} \setminus \ker \alpha_{1}|_p$.
    Definiamo $f : M \to \mathbb{R}$ puntuale come:
    \begin{equation*}
        f(p) = \frac{\alpha_{1}|_p(v_p)}{\alpha_{2}|_p(v_p)}
    \end{equation*}

    $f$ è indipendente dal vettore scelto nel sottospazio vettoriale 1-dimensionale $T_p \mathcal{M} \setminus \ker \alpha_{1}|_p$ perché qualsiasi vettore $U_p$ non-nullo nel sottospazio è risultato del riscalamento di $v_p$.
    Il fattore di riscalamento si semplifica nella definizione di $f$ grazie alla linearità di $\alpha_1, \alpha_2$.
    Per costruzione, $f$ è non-nulla e $\alpha_1 = f \alpha_2$.
    Dato che $\alpha_1, \alpha_2$ sono lisce, allora $f$ è liscia.
\end{proof}

Consideriamo il seguente esempio \cite{jiang_thermodynamics}: dotiamo $\mathbb{R}^{2n+1} \ni (x^1, y^1, \dots, x^n, y^n, z)$ della 1-forma:
\begin{equation*}
    \alpha = \odif{z} - \sum_{i=1}^{n} y^i \odif{x^i}
\end{equation*}
$(\mathbb{R}^{2n+1}, \alpha)$ è una varietà di contatto.

Nel caso $\mathbb{R}^3$, $\alpha = \odif{z} - y \odif{x}$ definisce la distribuzione di iperpiani:
\begin{equation*}
    \xi = \ker \alpha = \text{span}\{\partial_y, y \partial_z + \partial_x\}
\end{equation*}

Vediamo come ottenere la distribuzione valutando il $\ker \alpha$: cerchiamo tutti i vettori generici $V = A \partial_x + B \partial_y + C \partial_z \in T\mathbb{R}^3$ tali che $\alpha(V) = 0$.
(Nota: $\odif{x^i} \partial_{x_j} = \delta^i_j$)

\begin{align*}
    \alpha(V) & = (\odif{z} - y \odif{x}) (A \partial_x + B \partial_y + C \partial_z)                                                                                                                                                                                                        \\
              & = \cancelto{0}{A \odif{z} \partial_x} \quad - A y \: \cancelto{1}{\odif{x} \partial_x} \quad + \cancelto{0}{B \odif{z} \partial_y} \quad - \cancelto{0}{B y \odif{x} \partial_y} \quad + C \: \cancelto{1}{\odif{z} \partial_z} \quad - \cancelto{0}{C y \odif{x} \partial_z} \\
              & = C - Ay
\end{align*}

\begin{equation*}
    \alpha(V) = 0 \quad \Longrightarrow \quad C - Ay = 0 \quad \Longrightarrow \quad C = Ay
\end{equation*}
\begin{equation*}
    V_0 = A \partial_x + B \partial_y + A y \partial_z = B(\partial_y) + A(\partial_x + y \partial_z) \quad \Longrightarrow \quad \ker \alpha = \text{span}\{\partial_y, y \partial_z + \partial_x\} \quad \checkmark
\end{equation*}

Verifichiamo la condizione di non-integrabilità massimale:
\begin{align*}
    \odif{\alpha} & = - \sum_{i=1}^{n} \odif{(y^i \odif{x^i})}      \\
                  & = - \sum_{i=1}^{n} \odif{y^i} \wedge \odif{x^i}
\end{align*}
\begin{equation*}
    \Longrightarrow \alpha \wedge (\odif{\alpha})^n = -\left(\odif{z} - \sum_{i=1}^{n} y^i \odif{x^i}\right) \wedge \left(\sum_{j=1}^{n} \odif{y^j} \wedge \odif{x^j}\right)^n
\end{equation*}

\noindent \textbf{Nota:} Teorema multinomiale
\begin{equation*}
    \left(\sum_{i=1}^{r} X_i\right)^n = \sum_{\sum^r k_i = n} \binom{n}{k_1, \dots, k_r} \prod_{i=1}^{r} X_i^{k_i}
\end{equation*}

\begin{equation*}
    \left(\sum_{j=1}^{n} \odif{y^j} \wedge \odif{x^j}\right)^n = \sum_{\sum^r k_j = n} \binom{n}{k_1, \dots, k_r} \prod_{j=1}^{n} (\odif{y^j} \wedge \odif{x^j})^{k_j}
\end{equation*}
valutiamo il primo caso non banale ad es. $k_j = 2$: $(\odif{y^j} \wedge \odif{x^j})^2 = \odif{y^j} \wedge \odif{x^j} \wedge \odif{y^j} \wedge \odif{x^j} = - \odif{y^j} \wedge \odif{x^j} \wedge \odif{x^j} \wedge \odif{y^j} = 0$, grazie all'anticommutatività del prodotto wedge \(\omega \wedge \nu = - \nu \wedge \omega \Longrightarrow \omega \wedge \omega = 0\).

Quindi per $k_j \ge 2$ il termine $(\odif{y^j} \wedge \odif{x^j})^{k_j} = 0$.
Allora esiste una sola possibilità tale per cui non si annulli, cioè il caso $k_j = 1, \: \forall j $, da cui si vede $\sum^n_j k_j = n$.

\begin{align*}
    \left(\sum_{j=1}^{n} \odif{y^j} \wedge \odif{x^j}\right)^n & = \binom{n}{1,1,\dots,1} \prod_{j=1}^{n} (\odif{y^j} \wedge \odif{x^j})               \\
                                                               & = m! \, \odif{y^1} \wedge \odif{x^1} \wedge \dots \wedge \odif{y^n} \wedge \odif{x^n}
\end{align*}

\begin{align*}
    \Longrightarrow \alpha \wedge (\odif{\alpha})^n & = -\left(\odif{z} - \sum_{i=1}^{n} y^i \odif{x^i}\right) \wedge n! \bigwedge_{j=1}^{n} (\odif{y^j} \wedge \odif{x^j})                                                                          \\
                                                    & = -n! \, \odif{z} \wedge \bigwedge_{j=1}^{n} (\odif{y^j} \wedge \odif{x^j}) + \cancelto{0}{\sum_{i=1}^{n} x^i \odif{y^i} \wedge \left(\bigwedge_{j=1}^{n} \odif{y^j} \wedge \odif{x^j}\right)}
\end{align*}
(L'ultimo termine si annulla poiché $\odif{y^i} \wedge \odif{y^i} = 0$)

Quindi si ottiene:
\begin{equation*}
    \alpha \wedge (\odif{\alpha})^m = -m! \, \odif{z} \wedge \odif{x^1} \wedge \odif{y^1} \wedge \dots \wedge \odif{x^n} \wedge \odif{y^n}
\end{equation*}

in cui quest'ultima è una forma non-nulla su $\mathbb{R}^{2n+1}$. \checkmark

La forma \( \alpha = \odif{z} - \sum_{i=1}^{n} y^i \odif{x^i} \) è detta \textbf{struttura di contatto standard} su $\mathbb{R}^{2n+1}$.

\begin{definition}
    Sia ($\mathcal{M}, \xi$) una varietà di contatto; una sottovarietà immersa $L$ di ($\mathcal{M}, \xi$) è detta \textbf{isotropa} se $T_p L \subset \xi_p, \; \forall p \in L \subset \mathcal{M}$, cioè ogni spazio tangente è contenuto nella distribuzione.
\end{definition}

Possiamo ottenere un risultato molto importante sulle sottovarietà isotrope: $L$ è isotropa \(\iff i^* \alpha = 0\) dove \(i : L \longrightarrow \mathcal{M}\) è detta \emph{immersione canonica} cioè la funzione che mappa un punto di una sottovarietà in se stesso sulla varietà madre
\[
    p \in L \mapsto p \in \mathcal{M}, \quad i(p) = p
\]

Se prendiamo il \emph{pushforward} dell'immersione canonica, $\odif{i_p}$ prende un vettore $v \in T_p L$ e lo inietta formalmente in $T_p \mathcal{M}$, il che permette di passarlo in input ad $\alpha$ per riconoscerlo; il \emph{pullback} $i_p^*$ trasporta la 1-forma $\alpha : T\mathcal{M} \longrightarrow \mathbb{R}$ e ne restringe l'azione esclusivamente ad $L$. L'azione di queste due costruzioni si valuta tramite
\[
    (i^* \alpha )_p (v) = \alpha_p (\odif{i_p} v)
\]

Tramite queste ultime possiamo ottenere la giustificazione della condizione necessaria e sufficiente tra isotropia e l'annullamento del pullback della forma di contatto, infatti per ogni $v \in T_p L$,
\[
    (i^* \alpha )_p (v) = \alpha_p (\odif{i_p} v) = 0 \quad \iff \quad v \in \xi_p
\]
che vale perché $\xi = \ker \alpha$, quindi l'azione di $\alpha_p ( \odif{i_p}v) = 0 \quad \checkmark$

\section{Forme simplettiche}

La geometria di contatto è strettamente legata alla sua controparte su dimensioni pari, detta \emph{geometria simplettica}, poiché una forma di contatto induce un \emph{tensore simplettico} su ogni iperpiano di contatto

\begin{definition}
    Un 2-covettore $\omega$ che agisce su uno spazio vettoriale finito-dimensionale $\mathbb{V}$ è detto \textbf{non-degenere} o \textbf{tensore simplettico} o \textbf{forma bilineare simplettica} se
    \begin{equation*}
        \omega^n \neq 0 \tag*{Condizione di non-degenerazione}
    \end{equation*}
    $(\mathbb{V}, \omega)$ è detto \textbf{spazio vettoriale simplettico}.
\end{definition}

\begin{lemma}
    Sia $(\mathbb{V},\omega)$ uno spazio vettoriale simplettico. Allora i seguenti enunciati sono equivalenti:
    \begin{enumerate}
        \renewcommand{\theenumi}{\Roman{enumi}}
        \renewcommand{\labelenumi}{\theenumi)}
        \item $\omega^n \neq 0$ (non-degenerazione);
        \item la mappa lineare $\tilde{\omega}\colon \mathbb{V} \longrightarrow \mathbb{V}^*$, definita come
              $\tilde{\omega}(v) = \omega(v,\cdot)$, è invertibile;
        \item $\forall\, v \in V$ con $v \neq 0$, $\exists\, w \in V$ tale che $\omega(v,w) \neq 0$;
        \item esiste una base di $\mathbb{V}$ tale che la matrice $(\omega_{ij})$ che rappresenta $\omega$
              è non singolare.
    \end{enumerate}
\end{lemma}

\begin{proof}
    \noindent\textbf{II) $\Leftrightarrow$ III).}
    Per definizione,
    \[
        \tilde{\omega} \text{ è invertibile} \quad
        \iff \quad \tilde{\omega} \text{ è iniettiva} \quad
        \iff \quad \ker\tilde{\omega} = 0 ,
    \]
    % AGGIUNTA: motivazione dell'ultimo passaggio e dell'ipotesi di dimensione finita
    dove la prima equivalenza usa $\dim \mathbb{V} = \dim \mathbb{V}^* < \infty$. Infine $\ker\tilde{\omega}=0$ significa
    che per ogni $v\neq 0$ si ha $\omega(v,\cdot)\neq 0$, cioè III).

    \smallskip\noindent\textbf{II) $\Leftrightarrow$ IV).}
    Sia $\{e_i\}$ una base di $\mathbb{V}$ e $\{E^i\}$ la base duale di $\mathbb{V}^*$, cioè $E^j(e_i)=\delta^j_i$.
    In tale base la mappa lineare si scrive come
    \[
        \tilde{\omega}(e_i)(e_j) = \omega(e_i,e_j) = \omega_{ij} ,
        \qquad\text{cioè}\qquad
        \tilde{\omega}(e_i) = \omega_{ij}\,E^j
    \]
    (somma sugli indici ripetuti). Allora
    \begin{align*}
        \tilde{\omega} \text{ è invertibile}
         & \iff \exists\, \tilde{\omega}^{-1}\colon \mathbb{V}^*\to \mathbb{V} \text{ t.c. }
        \tilde{\omega}^{-1}\bigl(\tilde{\omega}(v)\bigr) = v \quad \forall\, v \in V         \\
         & \iff \tilde{\omega}^{-1}\bigl(\tilde{\omega}(e_i)\bigr) = e_i \quad \forall\, i   \\
         & \iff \omega_{ij}\,\tilde{\omega}^{-1}(E^j) = e_i \quad \forall\, i                \\
         & \iff \omega_{ij}\,E^k\bigl(\tilde{\omega}^{-1}(E^j)\bigr) = E^k(e_i) = \delta_i^k
        \quad \forall\, i,k                                                                  \\
         & \iff \exists\, \omega^{jk} \text{ t.c. } \omega_{ij}\,\omega^{jk} = \delta_i^k    \\
         & \iff \text{la matrice } (\omega_{ij}) \text{ che rappresenta } \omega
        \text{ è non singolare},
    \end{align*}
    dove si è posto $\omega^{jk} := E^k\bigl(\tilde{\omega}^{-1}(E^j)\bigr)$.

    \smallskip\noindent\textbf{III) $\Leftrightarrow$ I).}

    \noindent($\Rightarrow$) Costruiamo una base simplettica.
    \begin{definition}[Base simplettica]
        Sia $(\mathbb{V},\omega)$ uno spazio vettoriale simplettico. Una base
        $\{A_1,B_1,\dots,A_n,B_n\}$ di $\mathbb{V}$ è detta \emph{base simplettica} se
        \[
            \omega = \sum_{i=1}^n A^i \wedge B^i ,
        \]
        % AGGIUNTA: A^i, B^i sono la base duale
        dove $\{A^1,B^1,\dots,A^n,B^n\}$ è la base duale.
        Equivalentemente,
        \[
            \omega(A_i,A_j) = \omega(B_i,B_j) = 0
            \quad\text{e}\quad
            \omega(A_i,B_j) = -\,\omega(B_j,A_i) = \delta_{ij} .
        \]
    \end{definition}
    (Per la costruzione si veda~\cite{lee_manifolds}.)

    Avendo $\omega = \sum_{i=1}^n A^i \wedge B^i$, vale
    $\omega^n = \bigl(\sum_{i=1}^n A^i \wedge B^i\bigr)^n$.
    Posto $x_i = A^i\wedge B^i$ (sono $2$-forme, quindi commutano tra loro), ragioniamo come nel teorema multinomiale: esiste un'unica $n$-upla $(k_1,\dots,k_n)$ con $\sum_i k_i = n$ e $k_i \le 1$ per ogni $i$,
    cioè $(k_1,\dots,k_n) = (1,\dots,1)$. Dunque
    \[
        \omega^n = \Bigl(\sum_{i=1}^n A^i\wedge B^i\Bigr)^n
        = n!\, A^1\wedge B^1\wedge \dots \wedge A^n\wedge B^n \neq 0 .
    \]

    \medskip\noindent($\Leftarrow$) Per assurdo, supponiamo che III) non valga. Allora esiste $v\neq 0$
    tale che $\omega(v,\cdot)=0$, cioè $\ker\tilde{\omega}\neq 0$.
    Completiamo $v$ a una base $(v,E_2,\dots,E_{2n})$ di $\mathbb{V}$. Allora
    \[
        \omega^n(v,E_2,\dots,E_{2n}) = (\iota_v\omega^n)(E_2,\dots,E_{2n}) = 0 ,
    \]
    % AGGIUNTA: perché \iota_v \omega^n = 0 e perché basta valutare su una base
    perché $\iota_v\omega = \omega(v,\cdot) = 0$ implica
    $\iota_v\omega^n = n\,(\iota_v\omega)\wedge\omega^{n-1} = 0$.
    Essendo $\omega^n$ una $2n$-forma su uno spazio di dimensione $2n$, è determinata dal suo valore su una base, quindi $\omega^n = 0$, in contraddizione con I).
    Dunque vale III), cioè I) $\Longrightarrow$ III).
\end{proof}

Tornando alla definizione di tensore simplettico, possiamo notare una cosa importante: il fatto che dimensione di $\mathbb{V}$ sia pari può essere ottenuto dalla condizione di non-degenerazione.
Infatti un 2-covettore è un tensore covariante di rango 2 che agisce su $\mathbb{V} \times \mathbb{V}$, quindi è rappresentabile come una matrice $k \times k$.
È importante ricordare che una $k$-forma differenziale $\omega \in \Omega^k(\mathcal{M})$ assegna a ogni punto $p \in \mathcal{M}$ un tensore alternante $\omega_p \in \Lambda^k(T_p \mathcal{M})$.
$\omega_p$, per definizione, è una mappa multilineare che prende due vettori dello spazio tangente $T_p \mathcal{M}$ e restituisce un numero reale la 2-forma \(\omega_p : T_p \mathcal{M} \times T_p\mathcal{M} \longrightarrow \mathbb{R}\).
Scegliendo una base \( \{ \mathbf{e}_i \} \) per $T_p \mathcal{M}$, un generico vettore si scrive $v = v^i e_i$; calcolando l'azione della 2-forma su due vettori tangenti generici $v$ e $w$
\[
    \omega_p(v, w) = \omega(v^i e_i, w^j e_j) = v^i w^j \omega(e_i, e_j)
\]
I termini \(\omega_{ij} = \omega(e_i, e_j)\) sono i coefficienti della forma nella base scelta cioè tale che valga
\[
    \omega = \frac{1}{2} \omega_{ij} e^i \wedge e^j
\]
e formano la matrice quadrata richiesta.
La forma simplettica è alternante, quindi la sua rappresentazione in coordinate è una matrice $k \times k$ \emph{antisimmetrica}
\begin{align*}
    A^\intercal = - A \quad \Longrightarrow \quad \det A^\intercal & = \det (-A) \\
    \det A = (-1)^k \det A
\end{align*}
Valutiamo i due casi: \begin{itemize}
    \item $k$ dispari, vale che $\det A = - \det A \quad \Longrightarrow \det A = 0$; allora la matrice è singolare, ma il Lemma precedente impone la condizione di non-degenerazione (\(\omega^n \neq 0\)), da cui $\det \omega \neq 0$, quindi non è possibile;
    \item $k$ pari, $\det A = \det A$ \checkmark
\end{itemize}


Adesso possiamo estendere tali definizioni delle forme bilineari simplettiche su spazi vettoriali a una \emph{struttura simplettica} su varietà lisce.
\begin{definition}
    \cite{lee_manifolds} Sia $\mathcal{M}$ una varietà liscia, una 2-forma non-degenere su $\mathcal{M}$ è una 2-forma $\omega$ tale che $\omega_p$ è un 2-covettore non-degenere su $T_p \mathcal{M}$ per ogni $p \in \mathcal{M}$ (che rende $T_p \mathcal{M}$ uno spazio vettoriale simplettico).
    La 2-forma $\omega$ è detta \textbf{forma simplettica} o \textbf{struttura simplettica}.
\end{definition}

Si presentano due osservazioni: dato che i suoi spazi tangenti devono essere 2m-dimensionale, anche la varietà stessa sarà 2m-dimdimensionale; inoltre, su una varietà di contatto (2m+1)-dimensionale, la struttura di contatto associa un iperpiano 2m-dimdimensionale in ogni spazio tangente.
La struttura di contatto induce una forma bilineare simplettica su ciascun iperpiano.

\begin{lemma}
    Sia $(\mathcal{M}, \xi = \ker \alpha)$ una varietà di contatto, allora $(\xi_p, \odif{\alpha}|_{\xi_p})$ è un campo vettoriale simplettico dove $\odif{\alpha}|_{\xi_p}$ è definito come $\odif{\alpha}|_{\xi_p} : \xi_p \longrightarrow \xi_p^*$ tramite $v \mapsto \odif{\alpha}(v, \cdot)$
\end{lemma}
\begin{proof}
    La condizione di non-integrabilità massimale richiede che $(\odif{\alpha})^n \neq 0$ ovunque, quindi $\odif{\alpha}|_{\xi_p}$ soddisfa anch'essa la condizione di non-degenerazione per costruzione e sarà che $(\xi_p, \odif{\alpha}|_{\xi_p})$ è un campo vettoriale simplettico.
\end{proof}

Notare che lo spazio viene equipaggiato con la \emph{restrizione} della 2-forma $\odif{\alpha}$, $\odif{\alpha}|_{\xi_p}$ che agisce prendendo $v \in \xi_p$ e mandandolo in un covettore $\xi_p$ applicando $\odif{\alpha}(v, \cdot)$.

L'esistenza di una forma bilineare simplettica su uno spazio vettoriale fornisce una nozione generalizzata di ortogonalità, come per il prodotto interno, con l'aggiunta che i tensori simplettici sono prodotti antisimmetrici.

\begin{definition}
    Dato un sottospazio vettoriale $S$ di uno spazio vettoriale simplettico ($\mathbb{V}, \omega$), si definisce il \textbf{complemento simplettico}
    \[
        S^\perp = \{ v \in \mathbb{V} | \omega(v, w) = 0, \; \forall w \in S\}
    \]
\end{definition}

\section{Relazione tra le strutture simplettica e di contatto}
\begin{proposition}
    Sia \(
    (\mathcal{M}^{2n+1}, \xi = \ker \alpha)
    \) una varietà di contatto. Se $L$ è una sottovarietà isotropa di $\mathcal{M}$, allora $\dim L \leq n$.
\end{proposition}

\begin{definition}
    Se $\dim L = n$, diremo che $L$ è una sottovarietà \textbf{legendriana} di $\mathcal{M}$.
\end{definition}

Dimostriamo la proposizione:
\begin{proof}
    \cite{geiges_contact} Per ogni $p \in L$, poiché \[
        T_p L \subset \xi_p, \quad \Longrightarrow \quad i^* \alpha = 0
    \]
    inoltre \[
        \odif{(i^*\alpha)} = 0 \quad \Longrightarrow \quad i^* (\odif{\alpha}) = 0
    \] Quindi \( \forall u,v \in T_p L \), vale \( \odif{\alpha}|_{\xi|_p} (u,v) = 0 \). L'ultimo Lemma garantisce che, se dotato di $\odif{\alpha}|_{\xi_p}$, $\xi_p$ è un campo vettoriale simplettico, quindi possiamo definire il complemento simplettico di $T_p L$:
    \[
        (T_p L)^\perp = \{ v \in \xi_p | \odif{\alpha}|_{\xi_p} (u,v)= 0, \; \forall u \in T_p L \}
    \]
    Dato che \(  \odif{\alpha}|_{\xi_p} (u,v)= 0, \; \forall u, v \in T_p L  \), allora
    \[
        T_p L \subset (T_p L)^\perp \quad \Longrightarrow \quad \dim T_p L \leq \dim (T_p L)^\perp
    \]
    Costruiamo una mappa lineare \( \varphi : \xi_p \longrightarrow T_p^* L \), definita tramite \( \varphi (v)(w) = \odif{\alpha} (v, w)\) dove $v \in \xi_p$ e $w \in T_p L$.

    Dato che $\odif{\alpha}|_{\xi_p}$ è non-degenere, la mappa è un isomorfismo lineare suriettivo tale che
    \[
        \varphi(T_p L) \simeq T_p L \quad \Longrightarrow \quad \ker \varphi = (T_p L)^\perp
    \]

    Quindi abbiamo una funzione lineare tra spazi vettoriali per la quale il Teorema del Rango garantisce
    \begin{align*}
        \dim \xi_p & = \dim( \operatorname{Im} \varphi) + \dim( \ker \varphi) \\
        2 n        & = \dim T_p L + \dim (T_p L )^\perp
    \end{align*}
    Poiché \begin{align*}
                              & \dim T_p L \leq \dim (T_p L)^\perp \\
        \Longrightarrow \quad & \dim T_p L \leq n                  \\
        \Longrightarrow \quad & \dim L \leq n
    \end{align*}
\end{proof}

Questa proposizione e la sua dimostrazione mostrano due concetti chiave: per prima cosa, la distribuzione di iperpiani $\xi$ \emph{non} ammette \emph{varietà integrali}, definite \cite[Chapter 19]{lee_manifolds} come delle sottovarietà $\mathcal{N} \subseteq \mathcal{M}$ non-vuote immerse tali che
\[
    T_p \mathcal{N} = D_p, \; \forall p \in \mathcal{N}
\]
in cui $D$ è la distribuzione liscia $D \subseteq T \mathcal{M}$.
Quindi, la distribuzione di iperpiani $\xi$ non ammette una sottovarietà immersa $S$ che soddisfi $T_p S = \xi_p, \; \forall p \in S$.

Infatti le sottovarietà tangenti immerse hanno dimensione al più metà di quella della distribuzione.

Oltre a ciò, tale proposizione mostra come la geometria di contatto è strettamente legata alla geometria simplettica; si può utilizzare la nozione di varietà simplettica per creare una varietà di contatto.

Vediamo più nel dettaglio come effettuare questo salto.

Sia data una varietà differenziabile $\mathcal{M}$, $X$ campo vettoriale liscio su $\mathcal{M}$ e $F^X$ \emph{flusso} di $X$ nel senso \cite{latini_riemannian} che
\[
    F^X : \mathbb{R} \times \mathcal{M} \longrightarrow \mathcal{M} \; \text{ tale che } \; F^X(t, p) \coloneqq F^X_t(p) = \gamma_p (t)
\]
con $\gamma$ \emph{curva integrale} di X cioè
\[
    \gamma : I \subseteq \mathbb{R} \longrightarrow \mathcal{M} \; \text{ tale che } \; X(\gamma(t_0))=\dot\gamma (t_0)
\]
Allora \cite{lee_manifolds}, per un generico campo vettoriale $Y$ su $\mathcal{M}$ definiamo la \textbf{derivata di Lie} di $Y$ rispetto a $X$ come
\begin{equation}
    \left(\mathcal{L}_X Y\right)_p = \odv{}{t} \left.\odif{(F^X Y)}\right|_{t=0}
\end{equation}

\begin{definition}
    \cite{geiges_contact} Sia ($\mathcal{M}, \omega$) varietà simplettica (2n+2)-dimensionale, cioè $\omega$ è una 2-forma su $\mathcal{M}$ chiusa ($\odif{\omega} = 0$) e non-degenere ($\omega^{n+1} \neq 0$). Un campo vettoriale $X$ è detto \textbf{di Liouville} se
    \begin{equation}
        \mathcal{L}_X \omega = \omega
    \end{equation}
\end{definition}

Possiamo riscrivere la condizione di Liouville, tramite la \emph{formula di Cartan}
\begin{equation}
    \mathcal{L}_X = \operatorname{d} \circ \iota_X + \iota_X \circ \operatorname{d}
\end{equation}
dove bisogna prestare attenzione al fatto che $\iota_X$ è la \emph{contrazione} o \emph{prodotto interno}, non l'immersione canonica; da ciò otteniamo
\begin{align*}
    \mathcal{L}_X \omega                                                                  & = \omega \\
    \left( \operatorname{d} \circ \iota_X + \iota_X \circ \operatorname{d} \right) \omega & = \omega \\
    \odif{\iota_X \omega} + \iota_X \; (\cancelto{0}{\odif{\omega}} ) \quad               & = \omega \\
    \odif{\iota_X \omega}                                                                 & = \omega
\end{align*}

Allora la 1-forma $\alpha = \iota_X \omega$ definisce una \emph{forma di contatto} su qualsiasi $\xi$ in $\mathcal{M}$ trasversa a $X$.
Infatti, valutando la condizione di non-integrabilità massimale e sostituendo dentro
\[
    \alpha \wedge (\odif{\alpha})^n = \iota_X \omega \wedge \left(\odif{(\iota_X \omega)}\right)^n = \iota_X \omega \wedge \omega^n
\]
Diamo adesso un paio di proprietà che saranno utili in seguito.
\begin{property}
    Date due forme differenziali $\mu, \nu$ vale che
    \begin{equation}
        \mu \wedge \nu = (-1)^{\deg \mu \deg \nu} \nu \wedge \mu
    \end{equation}
    \begin{equation}
        \iota_X (\mu \wedge \nu) = (\iota_X \mu) \wedge \nu + (-1)^{\deg \mu} \mu \wedge (\iota_X \nu)
    \end{equation}
\end{property}
Prendiamo la forma \( \iota_X (\omega \wedge \omega^n) = (\iota_X \omega) \wedge \omega^n + \omega \wedge (\iota_X \omega^n) \), usando la prima proprietà appena espressa e unita al fatto che $\omega$ è una 2-forma (\( \deg \omega = 2 \)), otteniamo
\begin{align*}
    \omega \wedge (\iota_X \omega)   & = (\iota_X \omega) \wedge \omega   \\
    \omega^n \wedge (\iota_X \omega) & = (\iota_X \omega) \wedge \omega^n
\end{align*}
usiamo la seconda proprietà
\begin{align*}
    \iota_X (\omega^{n+1}) & = \iota_X (\omega \wedge \omega^n) =                                                                                                           \\
                           & = (\iota_X \omega) \wedge \omega^n + (-1)^{\deg \omega} \omega \wedge \iota_X \omega^n =                                                       \\
                           & = (\iota_X \omega) \wedge \omega^n + \omega \wedge \iota_X(\omega^n) =                                                                         \\
                           & = (\iota_X \omega) \wedge \omega^n + \omega \wedge \iota_X (\omega \wedge \omega^{n-1}) =                                                      \\
                           & = (\iota_X \omega) \wedge \omega^n + \omega \wedge \left[ (\iota_X \omega) \wedge \omega^{n-1} + \omega \wedge \iota_X (\omega^{n-2})\right] = \\
                           & = (\iota_X \omega) \wedge \omega^n + \omega \wedge (\iota_X \omega) \wedge \omega^{n-1} + \omega \wedge \omega \wedge \iota_X (\omega^{n-2})
\end{align*}
notiamo che nell'ultima espressione vale che
\[
    \omega \wedge (\iota_X \omega) \wedge \omega^{n-1} = (\iota_X \omega) \wedge \omega \wedge \omega^{n-1} = (\iota_X \omega) \wedge \omega^n
\]
Quindi
\begin{align*}
    \iota_X (\omega^{n+1})                                 & = 2(\iota_X \omega) \wedge \omega^n + \omega^2 \wedge \iota_X (\omega^{n-2}) = \\
    \text{\footnotesize{(iterando)}} \quad                 & = (n + 1) (\iota_X \omega \wedge \omega^n)                                     \\
    \Longrightarrow \qquad \iota_X \omega \wedge \omega^n  & = \frac{1}{n+1} \iota_X (\omega^{n+1})                                         \\
    \Longrightarrow \qquad \alpha \wedge (\odif{\alpha})^n & = \frac{1}{n+1} \iota_X (\omega^{n+1}) \neq 0 \checkmark                       \\
\end{align*}

\section{Applicazioni delle strutture}

Forniamo un metodo per confrontare due strutture e valutare se sono o meno la stessa:
\begin{definition}
    \cite{geiges_contact} Siano $\mathcal{M}_1, \alpha_1$ e $\mathcal{M}_2, \alpha_2$ due varietà di contatto: saranno dette \textbf{contattomorfe} se esiste un \emph{diffeomorfismo} $f : \mathcal{M}_1 \longrightarrow \mathcal{M}_2 $ tale che
    \[
        \odif{f}(\xi_1) = \xi_2 \quad \text{dove} \quad \odif{f} \coloneqq f_* : T\mathcal{M}_1 \longrightarrow T\mathcal{M}_2
    \]
    Se $ \xi_i = \ker \alpha $ per $ i = 1, 2 $, allora $f$ è detta \textbf{trasformazione di contatto} o \textbf{contattomorfismo}, e ciò equivale all'esistenza di una funzione liscia non-nulla $\lambda : \mathcal{M}_1 \longrightarrow \mathbb{R}$ tale che
    \[
        f^* \alpha_2 = \lambda \alpha_1 \quad \text{dove} \quad f^* : T^*\mathcal{M}_1 \longrightarrow T^*\mathcal{M}_2
    \]
    Infine, se $\lambda = 1$, diremo che il contattomorfismo è \textbf{esatto}.
\end{definition}

\subsection{1-jet e 1-grafici}
Introduciamo un paio di esempi che chiariscano l'utilità delle strutture simplettica e di contatto discusse fin qui.

Possiamo usare le funzioni lisce su una varietà per creare una varietà di contatto: poniamoci nello spazio euclideo per semplicità, ma potremmo generalizzarlo a varietà lisce generiche.
Sia $q$ un punto in $\mathbb{R}^n$ e $\varphi : \mathbb{R}^n \longrightarrow \mathbb{R}$ funzione liscia. Chiameremo \emph{1-jet} in $q$ di $\varphi$, etichettato come $J^1_q (\varphi)$, il polinomio di Taylor di $\varphi$ troncato al prim'ordine
\[
    J^1_q(\varphi) = \varphi(q) + \pdv{\varphi}{x^i} (q) x^i
\]
dove $ q = (x^1, \dots, x^n) \in \mathbb{R}^n $.

La nozione di 1-jet in $q$ di $\varphi$ è in relazione 1-1 con un punto di $\mathbb{R}^{2n + 1}$.
In questo modo si arriva a una nozione di \emph{struttura di contatto} naturale per lo spazio degli 1-jet.
Data una funzione liscia $\varphi$, il suo differenziale induce una sottovarietà legendriana: consideriamo l'\emph{1-grafico} di una funzione $\varphi$, detto anche \emph{prolungamento di ordine 1}
\[
    \Gamma(\varphi) = \left\{ \left( q, \varphi(q), \pdv{\varphi}{x}(q) \right) \right\} \subset J^1(\mathbb{R}^n, \mathbb{R})
\]
dove $ J^1(\mathbb{R}^n, \mathbb{R}) $ è l'insieme degli 1-grafici delle funzioni $ \mathbb{R}^n \longrightarrow \mathbb{R} $

L'1-forma $\alpha = \odif{y} - p \odif{x} $, dove $y$ è la coordinata $(n + 1)$-esima, corrispondente a $\varphi(q)$, e si annulla per pull back su tutti gli 1-grafici dalla mappa di immersione $i: \Gamma(\varphi) \longrightarrow J^1(\mathbb{R}^n, \mathbb{R})$
\[
    i^* \alpha = i^* (\odif{y} - p \odif{x}) = \odif{\varphi} - \pdv{\varphi}{x} \odif{x} = 0
\]
Quindi, gli 1-grafici di funzioni lisce sono sottovarietà legendriane.

\subsection{Fibrato cotangente proiettivizzato}

Data una qualsiasi varietà liscia $ (n + 1) $-dimensionale $B^{n + 1}$, si può creare una varietà di contatto tramite la seguente costruzione: un \emph{elemento di contatto} in un punto $p \in B$ è un iperpiano $H_p$ nello spazio tangente $T_p B$.
Si può pensare all'elemento di contatto come l'individuazione dell'insieme di $n$ derivate parziali di un 1-jet dell'esempio precedente.

La nozione di iperpiano $H_p$ è equivalente a una \emph{classe di equivalenza} di covettori $\alpha_p \in T_p^* B$ che annulla ogni vettore in $H_p$.
La relazione di equivalenza è ottenuta usando il fatto che due vettori linearmente dipendenti generano lo stesso iperpiano: ciò elimina la ridondanza che rendeva problematico associare a un iperpiano un insieme di covettori, indipendentemente dalla loro "magnitudine".

\begin{remark}
    Ricordiamo che, come espresso precedentemente da \cite[Capitolo 2, Foundations of Special Relativity]{misner-wheeler-thorne_gravitation} Misner, Thorne, Wheeler, la norma di un covettore rappresenta la densità del fascio di piani ortogonali al vettore tangente che formano la nozione intuitiva di covettori.
\end{remark}

Quindi, gli elementi di contatto in $p \in B$ sono anche classi di equivalenza di covettori in $T_p^* B$.
Lo spazio di tutti gli elementi di contatto al punto $p$, per ogni $p \in B$, può essere rappresentato come il \emph{fibrato cotangente proiettivizzato} $\mathbb{P}T_p^* B$.
Tale spazio ha struttura di fibrato dove una fibra esemplificativa in $p$ è lo spazio dei raggi in $T_p^* B$.

\begin{remark}
    \begin{enumerate}
        \item un raggio è una retta vettoriale (cioè un sottospazio 1-dimensionale)
        \item la proiettivizzazione è la formalizzazione più rigorosa di ciò che abbiamo eseguito tramite l'associazione a una "direzione" di covettori tutti i covettori che giacciono su di essa
    \end{enumerate}
\end{remark}
Scegliendo una rappresentazione, un punto di $\mathbb{P}T_p^* B$ è scrivibile come
\[
    (u, x^1, \dots, x^n, p^1, \dots, p^n)
\]
dove $u, x^1, \dots,x^n$ sono coordinate per B e $p^1, \dots, p^n$ sono le coordinate nella base per i covettori.
Riprendendo l'esempio degli 1-jet, $u$ sono i possibili valori di una funzione di $x^1, \dots, x^n$ a valori reali e $p^1, \dots, p^n$ i possibili valori delle derivate parziali della funzione rispetto alle $x^i$ in quel punto di $B$.

\subsection{Meccanica hamiltoniana}

Queste strutture descritte fino a qui emergono in vari modi in fisica: uno dei modi più noti è la formulazione lagrangiana della meccanica classica.
Assumiamo che i sistemi siano descritti da un'azione $S$, definita come
\[
    S[\gamma(t)] = \int_\gamma \mathcal{L} (q, \dot q, t) \odif{t}
\]
dove $S$ è un funzionale di un cammino $\gamma$ nello spazio delle fasi e $q$ sono le coordinate generalizzate; la soluzione classica delle equazioni del moto è il percorso $\gamma(t)$ che estremizza l'azione $S$ ed è ottenibile risolvendo le equazioni di Eulero-Lagrange.
La formulazione hamiltoniana è collegata a quella lagrangiana tramite la trasformata di Legendre: sia $H(q, p, t)$ l'hamiltoniana del sistema, dove $p$ sono i momenti coniugati $\frac{\mathcal{L}}{\dot q}$, allora, data la relativa lagrangiana $\mathcal{L}(q, \dot q, t)$, scriveremo
\[
    H(q, p, t) = p \dot q - \mathcal{L}(q, \dot q, t) \qquad \text{quindi} \qquad S = \int_\gamma (p \dot q - H) \odif{t}
\]

Supponiamo che il pullback agisca sulle forme differenziali di $\mathcal{M}$ riportandole su $\mathbb{R}$, dove $\mathbb{R}$ è la varietà di partenza del tempo con vettori tangenti $\odv{}{t}$ e $\mathcal{M}$, la varietà di arrivo, è lo spazio delle fasi hamiltoniano $\{ (q, p, t) \}$: consideriamo $\gamma : \mathbb{R} \longrightarrow \mathcal{M}$ che parametrizza le traiettorie in $\mathcal{M}$ tramite la coordinata temporale.
Quindi
\[
    (\gamma^* \alpha) (v) = \alpha (\odif{\gamma}(v))
\]

Valutiamo
\begin{equation}
    \alpha = p \odif{q} - H \odif{t} \tag*{\textbf{Forma di Poincaré-Cartan}}
\end{equation}
Poniamo $v = \odv{}{t}$, quindi il vettore
\begin{align*}
    \odif{\gamma} \left( \odv{}{t} \right) & = \odv{q}{t} \pdv{}{q} + \odv{p}{t} \pdv{}{p} + \odv{t}{t} \pdv{}{t} = \\
                                           & = \dot q \pdv{}{q} + \dot p \pdv{}{p} + \pdv{}{t}
\end{align*}
da cui, $\alpha (\odif{\gamma}(v))$ :
\begin{align*}
    \alpha\left(\odif{\gamma}\left(\odv{}{t}\right)\right) & = (p \odif{q} - H \odif{t}) \left( \dot q \pdv{}{q} + \dot p \pdv{}{p} + \pdv{}{t} \right) =                                                                            \\
                                                           & = p \dot q \odif{q} \pdv{}{q} - H \dot q \odif{t} \pdv{}{q} + \dot p p \odif{q} \pdv{}{p} - H \dot p \odif{t} \pdv{}{p} + p \odif{q} \pdv{}{t} - H \odif{t} \pdv{}{t} =
\end{align*}

\begin{remark}
    Sappiamo che vale $\odif{x^i}(\partial x_j) = \delta^i_j$
\end{remark}
Otteniamo
\begin{align*}
    \alpha \left( \odif{\gamma} \left( \odv{}{t} \right) \right)                   & = p \dot q - H                                     \\
    \Longrightarrow \qquad \left( \gamma^* \alpha \right) \left( \odv{}{t} \right) & = p \dot q - H                                     \\
                                                                                   & = (p \dot q - H) \odif{t} \left( \odv{}{t} \right) \\
    \Longrightarrow \qquad \left( \gamma^* \alpha \right)                          & = (p \dot q - H) \odif{t}
\end{align*}
sostituendo, si ottiene infine:
\begin{equation}
    S = \int_\gamma \gamma^* \alpha
\end{equation}